\originalchapter{环、理想与商环}
\originalsection{两种运算}
\begin{definition}
本书通常讨论有乘法单位元且 $1\ne0$ 的环; 为使商环和同构定理包含全理想情形, 也允许零环作为退化例外, 其 $1=0$. 域和整环始终要求 $1\ne0$. 加法使 $R$ 成交换群, 乘法结合且满足左右分配律. 不额外说“交换”时允许乘法不交换. 环同态保持加法、乘法和单位元. 子环包含同一单位元. 非零元素若有乘法逆元称为单位. 交换环若无非零零因子称为整环; 若每个非零元素可逆称为域.
\end{definition}
整数环、域、多项式环和矩阵环都是例子. $\Z/n\Z$ 中 $\bar a$ 可逆当且仅当 $\gcd(a,n)=1$: 贝祖等式给出逆元, 反之 $ab\equiv1$ 说明任何公因子整除1. 因此 $\Z/n\Z$ 是域当且仅当 $n$ 素数. 非零有限整环总是域: 对 $a\ne0$, 映射 $x\mapsto ax$ 单射, 有限性给出满射, 因而1在其像中.

$0a=(0+0)a=0a+0a$, 加法消去得 $0a=0$, 也有 $a0=0$. 再由 $a+(-a)=0$ 得 $(-a)b=-(ab)$, 最后 $(-a)(-b)=ab$. 这些是分配律的结果, 不能另作乘法公理.
\originalsection{理想与同构定理}
\begin{definition}
环 $R$ 的双边理想 $I$ 是加法子群, 满足 $rI\subseteq I$, $Ir\subseteq I$ 对每个 $r\in R$ 成立. 交换环中只需写 $rI\subseteq I$. 由 $a$ 生成的主理想记 $(a)$; 交换环中 $(a)=\{ra:r\in R\}$.
\end{definition}
理想通常不含1; 若含1则是全环. 环同态的核是理想, 因 $\varphi(ri)=\varphi(r)\varphi(i)=0$, 右乘同理. 对理想 $I$, 在加法陪集上定义 $(a+I)(b+I)=ab+I$: 改变代表元时增量为 $ai+jb+ji\in I$, 故良定义. 此商环 $R/I$ 的单位元为 $1+I$; 真理想保证 $1+I\ne0$. 若 $I=R$, 商为零环, 本书不把它当作域.
\begin{theorem}
对单位保持环同态 $\varphi:R\to S$, 有 $R/\ker\varphi\cong\im\varphi$. 包含 $I$ 的理想与 $R/I$ 的理想由像和原像一一对应. 若 $I\subseteq J$, 则 $(R/I)/(J/I)\cong R/J$.
\end{theorem}
\begin{proof}
映射 $a+\ker\varphi\mapsto\varphi(a)$ 与群情形相同地良定义且双射, 又保持乘法和1. 理想的原像满足加法封闭及两边吸收; 若原理想包含核, 其像在满射的像环中也有吸收性. 原像与像互逆的证明和群对应定理相同. 最后映射 $R/I\to R/J$ 的核为 $J/I$, 用第一结论.
\end{proof}
\originalsection{素理想与极大理想}
\begin{definition}
以下 $R$ 交换. 真理想 $P$ 若 $ab\in P$ 必有 $a\in P$ 或 $b\in P$, 称素理想. 真理想 $M$ 若它和 $R$ 之间无其他理想, 称极大理想.
\end{definition}
\begin{theorem}\label{thm:maximal}
$R/P$ 为整环当且仅当 $P$ 为素理想; $R/M$ 为域当且仅当 $M$ 为极大理想. 因此极大理想必为素理想.
\end{theorem}
\begin{proof}
商环无零因子恰为 $ab\in P$ 迫使一个因子在 $P$ 中. 若 $M$ 极大, 对 $a\notin M$ 有 $M+(a)=R$, 所以存在 $m\in M,r\in R$ 使 $m+ra=1$, 从而 $a+M$ 可逆. 反之若商为域, 对理想 $J\supsetneq M$ 取 $a\in J\setminus M$, 其像可逆, 故 $J/M$ 含1, 即 $J=R$.
\end{proof}
\originalsection{中国剩余定理}
理想 $I,J$ 的和由 $i+j$ 组成, 交为普通集合交, 积 $IJ$ 由有限和 $\sum i_kj_k$ 组成; 逐项分配可证它们都是理想. 总有 $IJ\subseteq I\cap J$, 但一般不相等.
\begin{theorem}[中国剩余]\label{thm:crt}
交换环中若 $I+J=R$, 则 $I\cap J=IJ$, 且 $R/IJ\cong R/I\times R/J$.
\end{theorem}
\begin{proof}
取 $i+j=1$, $i\in I,j\in J$. 对 $x\in I\cap J$, 写 $x=xi+xj\in IJ$, 故交等于积. 映射 $r\mapsto(r+I,r+J)$ 的核是交. 给定 $(a+I,b+J)$, 元素 $aj+bi$ 映到它, 因 $j\equiv1$ 模 $I$ 且 $i\equiv1$ 模 $J$, 所以满射. 同构定理给出结论.
\end{proof}
\begin{example}构造 $\Z/15\Z$ 的非平凡幂等元, 并说明有限环的单位与零因子.\end{example}
\begin{solution}
$\Z/15\Z\cong\F_3\times\F_5$. 幂等元在两个域中各只能0或1, 所以共有4个: $0,1,6,10$. $6^2\equiv6$, $10^2\equiv10$ 模15. 互素于15的元素为单位, 其余非零元素为零因子, 如 $\bar3\bar5=0$.
\end{solution}
\originalsection{带解答练习}
\begin{exercise}证明 $\Z[x]$ 的 $(x)$ 为素理想却不极大, 而 $(2,x)$ 为极大理想.\end{exercise}
\begin{solution}
代入0的满射 $\Z[x]\to\Z$ 的核是 $(x)$, 商为整环而非域. 模2再代入0的满射到 $\F_2$ 核为 $(2,x)$, 商为域.
\end{solution}
\begin{exercise}说明直积环 $R\times S$ 一般不满足任意两环到第三环的余积泛性质 (原附注的坐标注入采用非单位保持约定).\end{exercise}
\begin{solution}
单位保持约定下 $r\mapsto(r,0)$ 不是单位保持映射. 即便允许非单位保持同态, 对 $R=S=T=\Z$ 的两个恒等映射, 想要的加法映射必为 $(a,b)\mapsto a+b$, 但 $(1,0)(0,1)=0$ 的像为0, 两个像的乘积为1. 因而不保持乘法. 原作业8的相关附注须改正; 阿贝尔群和模的有限直和确实与直积相同, 环的情况不能照搬.
\end{solution}
